\section{Objetivos}

\subsection{Objetivo general}

Diseñar, implementar y evaluar un sistema reproducible de detección no
supervisada de anomalías en radiografías de tórax, comparando un AE, un VAE y
una alternativa adversarial GAN bajo un protocolo común.

\subsection{Objetivos específicos y criterios de logro}

\begin{enumerate}
  \item \textbf{Revisar el estado del arte} de detección de anomalías basada en
  reconstrucción, identificando supuestos, funciones de pérdida, puntuaciones y
  riesgos. Se considera cumplido al justificar las tres familias y el protocolo
  a partir de bibliografía primaria.

  \item \textbf{Obtener y preparar un dataset médico adecuado.} Debe contener
  radiografías normales y anómalas, entrenamiento normal, validación y prueba.
  Se verifica mediante conteos, lectura de imágenes, dimensiones y control de
  duplicados exactos entre particiones.

  \item \textbf{Implementar tres modelos comparables}: AE convolucional, VAE y
  GANomaly. Deben compartir resolución, dimensión latente, optimizador,
  presupuesto de épocas y semillas, y producir una puntuación por imagen.

  \item \textbf{Definir una evaluación sin fuga de información.} Los pesos se
  ajustan solo con entrenamiento normal; el checkpoint se selecciona con
  validación normal; las etiquetas completas de validación fijan el umbral; el
  test se consulta una vez con el método congelado.

  \item \textbf{Comparar efectividad y coste}. La métrica principal es AUROC de
  test; se reportan AUPRC, exactitud balanceada, sensibilidad, especificidad,
  F1, número de parámetros y tiempo. La afirmación de superioridad requiere la
  media de tres semillas y dispersión, no una ejecución aislada.

  \item \textbf{Analizar limitaciones y reproducibilidad}. Se documentan
  desequilibrio, sesgos de adquisición, pérdida por redimensionado, significado
  limitado de la etiqueta y ausencia de validación clínica externa.
\end{enumerate}

Quedan fuera del alcance mínimo la localización de lesiones, la clasificación
por patología, el uso asistencial y la comparación exhaustiva con métodos
preentrenados modernos. Son extensiones válidas, pero diluirían la pregunta
central y el presupuesto de un TFM individual de 18 ECTS.
